\begin{proof}[Proof of \cref{obj:lem:common-marginal-canonical-certificate-family}]
Set
\[
C_{\mathrm{fixed}}=65536,
\qquad
c_{\mathrm{risk}}=\frac1{65536}.
\]
These constants are positive.  Fix \(\epsilon\in(0,1/2)\).  By \cref{obj:lem:common-marginal-uniform-intensity}, choose a common-marginal calibration and constants \(C_\epsilon,\rho_\epsilon\) with
\[
C_\epsilon>0,\qquad \rho_\epsilon\in(0,1),
\]
together with the uniform construction and closure bounds supplied there.  Define
\[
c_{\mathrm{small}}=\frac{1}{65536\,C_\epsilon},
\qquad
c_{\mathrm{floor}}=\frac1{65536}.
\]
Then
\[
c_{\mathrm{small}}>0,\qquad c_{\mathrm{floor}}>0,\qquad
c_{\mathrm{floor}}\le c_{\mathrm{risk}},
\]
and \cref{obj:lem:common-marginal-uniform-intensity} supplies positive
\(b_0,b_\epsilon,\gamma_\epsilon,q_\epsilon,c_\epsilon,C_d\), with
\(c_\epsilon\le 1/4\).  We retain that statement's notation: \(q_\epsilon\) is the
rare-count coefficient, \(c_\epsilon\) is the dual-gap constant, and \(C_d\) is
the dimension constant. % lean: common_marginal_canonical_certificate_family

First consider the fixed-sample transfer assertion.  Let \(n,m,d,H\) satisfy the hypotheses in that assertion, and write
\[
\Delta
=
\left|\theta_1(H)-\theta_0(H)\right|.
\]
By the hypotheses of the fixed-sample assertion, the recipe carried by \(H\) comes with an exact Poisson-prior witness at
\[
u=C_{\mathrm{fixed}}n,
\qquad
v=C_{\mathrm{fixed}}(n+m),
\]
with constants \(C_\epsilon,\rho_\epsilon\).  The same hypotheses give total variation at most \(1/8\) for its prior-predictive laws, and the assumed variance absorption is
\[
kBa\le c_{\mathrm{small}}\Delta^2
=
\frac{1}{65536\,C_\epsilon}\Delta^2 .
\]
Applying the fixed-sample recipe transfer bound \cref{obj:lem:common-marginal-recipe-transfer} at these intensities gives
\[
\frac1{65536}\Delta^2
\le
R_\epsilon(n,m,d).
\]
Since \(c_{\mathrm{risk}}=1/65536\), this is exactly the fixed-sample transfer conclusion. % lean: fixed_sample_common_marginal_transfer_rule

It remains to construct the canonical handle for arbitrary \(n,m,d\) with \(n\ge1\) and \(d\ge2\).  Write
\[
N=n+m,\qquad \ell_n=\log(en),
\qquad
D(n,m,d)=\min\left\{1,\frac{d^2}{N^2\ell_n^2}\right\}.
\]
Let \(r_\epsilon\) denote the calibration regime constant, \(D_\epsilon\) the calibration degree constant, and \(q_\epsilon\) the calibration rare-count constant.  The proof splits according to whether the parametric regime holds:
\[
D(n,m,d)\le \frac{r_\epsilon}{n}.
\]
In both regimes set
\[
L=\left\lceil D_\epsilon\ell_n\right\rceil,\qquad
B=\frac{b_0L}{N},\qquad
a=\frac{\gamma_\epsilon B}{L^2},
\]
and
\[
k=\min\left\{d-1,\left\lfloor q_\epsilon NL\right\rfloor\right\}.
\]
The calibration closure in \cref{obj:lem:common-marginal-uniform-intensity} gives
\[
L\ge2
\]
and
\[
2C_{\mathrm{fixed}}b_0=131072\,b_0\le b_\epsilon .
\]
Moreover, since \(n\le N\),
\[
(C_{\mathrm{fixed}}n+C_{\mathrm{fixed}}N)B
\le
(131072\,N)\frac{b_0L}{N}
=
131072\,b_0L
\le b_\epsilon L.
\]
Thus the degree and bandwidth conclusions hold for either branch. % lean: common_marginal_canonical_certificate_family

Suppose first that
\[
D(n,m,d)\le \frac{r_\epsilon}{n}.
\]
The local testing gap used by the parametric handle is admissible because
\[
\frac{c_\epsilon}{\sqrt n}\le c_\epsilon\le \frac14,
\]
where \(\sqrt n\ge1\).  Let \(H\) be the parametric common-marginal prior handle with the above values of \(L,B,a,k\).  Its branch condition makes the finite-atom predicate false in this case.  Consequently all conclusions whose hypotheses include the finite-atom predicate are vacuous.  The common treatment-covariate marginal identity is the handle identity
\[
\mathcal L(P_{XA}\mid \Pi_0(H))
=
\mathcal L(P_{XA}\mid \Pi_1(H)).
\]
The risk floor follows from the parametric two-point floor:
\[
\frac{c_{\mathrm{floor}}}{n}
=
\frac{1/65536}{n}
\le
\frac{1/100}{n}
\le
R_\epsilon(n,m,d).
\]
Finally, the calibration closure gives
\[
C_d\,r_\epsilon\le \frac1{65536}=c_{\mathrm{floor}},
\]
and hence, using the parametric-regime inequality,
\[
C_d\,D(n,m,d)
\le
C_d\,\frac{r_\epsilon}{n}
\le
\frac{c_{\mathrm{floor}}}{n}.
\]
This proves all required canonical-handle properties on the parametric branch. % lean: common_marginal_canonical_certificate_family

Now suppose that the parametric regime fails:
\[
D(n,m,d)>\frac{r_\epsilon}{n}.
\]
With \(L,B,a,k\) as above, the quantities \(B\) and \(a\) are positive.  From
\[
k=\min\left\{d-1,\left\lfloor q_\epsilon NL\right\rfloor\right\}
\]
and the elementary inequality \(\lfloor q_\epsilon NL\rfloor\le q_\epsilon NL\), valid since \(q_\epsilon NL\ge0\), we have
\[
k\le q_\epsilon NL.
\]
Therefore
\[
ka
\le
q_\epsilon NL\cdot
\frac{\gamma_\epsilon b_0}{NL}
=
q_\epsilon\gamma_\epsilon b_0
\le c_\epsilon,
\]
where the last inequality is the rare-scale calibration inequality.  Hence the rare shift is small. % lean: common_marginal_canonical_certificate_family

Apply the finite-atom part of \cref{obj:lem:common-marginal-uniform-intensity} with
\[
u=65536\,n,\qquad v=65536\,N.
\]
The preceding bandwidth bound, positivity of \(B\), the identity \(a=\gamma_\epsilon B/L^2\), the rare-shift bound \(ka\le c_\epsilon\), and \(k+1\le d\) produce a recipe \(R\) at these intensities.  The bounds supplied for this recipe include the finite-atom construction flag, support in \(\mathcal M_{d,\epsilon}\), equality of the induced treatment-covariate marginal table, the raw total-mass mean-one identity, the raw-center separation, the raw total-mass and target variance bounds, and the raw mixture total-variation bound
\[
\operatorname{TV}
\bigl(\mathsf P_{0,R}^{\mathrm{Pois}},\mathsf P_{1,R}^{\mathrm{Pois}}\bigr)
\le
C_\epsilon u k a \rho_\epsilon^L.
\]
Together with \(n\ge1\), \(L\ge2\), \(d\ge2\), the nonnegativity of \(u,v\), the positivity of \(u+v\), the bandwidth and shift identities, and the rare-count fit \(k+1\le d\), these same bounds form the exact Poisson-prior witness used below.  Let \(H\) be the corresponding handle.  Its branch condition identifies this recipe with the finite-atom branch.  The common marginal identity again gives
\[
\mathcal L(P_{XA}\mid \Pi_0(H))
=
\mathcal L(P_{XA}\mid \Pi_1(H)).
\]
The recipe bounds also give
\[
\left|\theta_1(H)-\theta_0(H)\right|
\ge
c_\epsilon ka
=
\Delta_H,
\]
the raw total mass has mean one, the raw total mass and raw targets are square-integrable with variances at most \(C_\epsilon kBa\), and the raw mixture total variation is bounded by
\[
C_\epsilon u k a\rho_\epsilon^L.
\]
The calibration decay inequality and the bound \(ka\le c_\epsilon\) imply
\[
\operatorname{TV}
\bigl(\mathsf P_{0,H}^{\mathrm{Pois}},\mathsf P_{1,H}^{\mathrm{Pois}}\bigr)
\le
C_\epsilon(65536\,n)ka\rho_\epsilon^L
\le
C_\epsilon(65536\,n)c_\epsilon\rho_\epsilon^L
\le
\frac18 .
\]
Thus the prior-predictive total variation condition required by the transfer rule holds. % lean: common_marginal_canonical_certificate_family

It remains to record the two absorptions used by the handle.  The counting estimate used in the finite-atom branch is
\[
\frac{L^2}{c_{\mathrm{small}}c_\epsilon^2\gamma_\epsilon}\le k .
\]
Indeed, write \(N=n+m\) and \(\ell_n=\log(en)\).  Since the branch condition is
\[
D(n,m,d)>\frac{r_\epsilon}{n},
\]
the minimum in \(D(n,m,d)\) forces
\[
\frac{r_\epsilon}{n}
<
\frac{d^2}{N^2\ell_n^2}.
\]
Put
\[
K=c_{\mathrm{small}}c_\epsilon^2\gamma_\epsilon>0.
\]
The calibration inequality
\[
\left(
\frac{4(D_\epsilon+1)^2}{K}
\right)^2
\le r_\epsilon
\]
implies
\[
\frac{4(D_\epsilon+1)^2}{K}\le\sqrt{r_\epsilon}.
\]
Meanwhile, the failed-regime inequality above gives
\[
\frac{\sqrt{r_\epsilon}\,N\ell_n}{\sqrt n}<d.
\]
Using the ceiling bound \(L\le(D_\epsilon+1)\ell_n\) and
\[
\ell_n\sqrt n\le 2n\le 2N,
\]
we obtain the complete cross-multiplication
\[
\begin{aligned}
\frac{2L^2}{K}
&\le
\frac{2(D_\epsilon+1)^2\ell_n^2}{K} \\
&=
\left(\frac{4(D_\epsilon+1)^2}{K}\right)
\frac{N\ell_n}{\sqrt n}
\frac{\ell_n\sqrt n}{2N} \\
&\le
\frac{\sqrt{r_\epsilon}\,N\ell_n}{\sqrt n}
<d.
\end{aligned}
\]
Equivalently, the proof has compared the nonnegative quantities through the strict square inequality
\[
\left(\frac{4(D_\epsilon+1)^2}{K}N\ell_n\right)^2
<
(d\sqrt n)^2,
\]
and then divided by the positive factor \(\sqrt n\).  Weakening the strict bound gives
\[
\frac{2L^2}{K}
\le d.
\]
Now split according to the actual comparison defining the minimum in \(k\).  If
\[
d-1\le \left\lfloor q_\epsilon NL\right\rfloor,
\]
then
\[
k=d-1.
\]
Since \(d\ge2\),
\[
\frac d2\le d-1,
\]
and hence
\[
\frac{L^2}{K}
\le \frac d2\le d-1=k .
\]
If instead
\[
\left\lfloor q_\epsilon NL\right\rfloor<d-1,
\]
then
\[
k=\left\lfloor q_\epsilon NL\right\rfloor .
\]
The calibration gives
\[
\frac{4(D_\epsilon+1)}{K}
\le q_\epsilon,
\]
and, since \(q_\epsilon NL\ge1\), the elementary floor bound \(\lfloor q_\epsilon NL\rfloor\ge q_\epsilon NL/2\) gives
\[
k=\bigl\lfloor q_\epsilon NL\bigr\rfloor
\ge \frac{q_\epsilon NL}{2}
\ge
\frac{2(D_\epsilon+1)NL}{K}.
\]
Using \(L\le(D_\epsilon+1)\ell_n\) and \(\ell_n\le N\), this is at least
\[
\frac{L^2}{K}.
\]
Thus the displayed counting estimate holds in both subcases.  Recall here that \(c_\epsilon\) denotes the calibration dual-gap constant; the rare-count coefficient is \(q_\epsilon\).  Using \(a=\gamma_\epsilon B/L^2\), the counting estimate yields
\[
kBa
=
\left(\frac{k a^2}{\gamma_\epsilon}\right)L^2
\le
\frac{k a^2}{\gamma_\epsilon}
\bigl(kc_{\mathrm{small}}c_\epsilon^2\gamma_\epsilon\bigr)
=
c_{\mathrm{small}}(c_\epsilon ka)^2
=
c_{\mathrm{small}}\Delta_H^2.
\]
The recipe separation bound gives
\[
\Delta_H=c_\epsilon ka
\le
\left|\theta_1(H)-\theta_0(H)\right|.
\]
Consequently
\[
kBa
\le
c_{\mathrm{small}}
\left|\theta_1(H)-\theta_0(H)\right|^2 .
\]
The fixed-sample transfer already proved above therefore gives
\[
c_{\mathrm{floor}}
\left|\theta_1(H)-\theta_0(H)\right|^2
\le
R_\epsilon(n,m,d),
\]
and hence
\[
c_{\mathrm{floor}}\Delta_H^2\le R_\epsilon(n,m,d).
\]
As before, the parametric floor gives
\[
\frac{c_{\mathrm{floor}}}{n}\le R_\epsilon(n,m,d).
\]
Finally we verify the dimension absorption for the above \(L,k,a\).  Since
\[
a=\frac{\gamma_\epsilon B}{L^2}
=\frac{\gamma_\epsilon b_0}{NL},
\]
it is enough to compare \(C_dD(n,m,d)\) with
\[
\left(c_\epsilon k\frac{\gamma_\epsilon b_0}{NL}\right)^2.
\]
Split according to which term realizes
\[
k=\min\left\{d-1,\left\lfloor q_\epsilon NL\right\rfloor\right\}.
\]
On the cap-active subcase, \(k=d-1\).  Because \(d\ge2\),
\[
\frac d2\le d-1=k.
\]
The definition of \(D(n,m,d)\) and the dimension calibration give
\[
\begin{aligned}
C_dD(n,m,d)
&\le C_d\frac{d^2}{N^2\ell_n^2} \\
&\le
\left(\frac{c_\epsilon\gamma_\epsilon b_0}{2(D_\epsilon+1)}\right)^2
\frac{d^2}{N^2\ell_n^2}.
\end{aligned}
\]
All factors here are nonnegative.  The two bounds
\(d\le2k\) and \(L\le(D_\epsilon+1)\ell_n\) therefore give, after
cross-multiplication by the positive denominators,
\[
0\le
\frac{c_\epsilon\gamma_\epsilon b_0d}
     {2(D_\epsilon+1)N\ell_n}
\le
c_\epsilon k\frac{\gamma_\epsilon b_0}{NL}
=c_\epsilon ka.
\]
Squaring this nonnegative comparison proves
\[
C_dD(n,m,d)\le(c_\epsilon ka)^2
\]
on the cap-active subcase.
On the floor-active subcase, \(k=\lfloor q_\epsilon NL\rfloor\).  The calibration includes the rare-count identity
\[
q_\epsilon=\frac{c_\epsilon}{2\gamma_\epsilon b_0}.
\]
The calibration has \(q_\epsilon\ge4\), while \(N\ge1\) and \(L\ge2\), so
\(q_\epsilon NL\ge1\).  The elementary floor bound therefore gives
\[
k\ge \frac{q_\epsilon NL}{2}.
\]
Thus, with every factor nonnegative,
\[
ka
\ge
\frac{q_\epsilon NL}{2}\cdot\frac{\gamma_\epsilon b_0}{NL}
=\frac{c_\epsilon}{4}.
\]
Since \(D(n,m,d)\le1\), the calibration gives
\(C_d\le c_\epsilon^4/16\), and squaring preserves the preceding comparison,
\[
C_dD(n,m,d)
\le C_d
\le \frac{c_\epsilon^4}{16}
=\left(c_\epsilon\frac{c_\epsilon}{4}\right)^2
\le (c_\epsilon ka)^2.
\]
Thus in both subcases
\[
C_d\,D(n,m,d)\le (c_\epsilon ka)^2=\Delta_H^2.
\] % lean: commonMarginal_dimension_absorption
This proves every required canonical-handle property on the finite-atom branch, and completes the construction for all \(n,m,d\). % lean: common_marginal_canonical_certificate_family
\end{proof}
